\subsection{Exact marking transport through a finite normal and mirror core}

\begin{result}{A.4.5}{Proved}{The marking ambiguity has an explicit finite refinement}
Let $B=C_2*C_3=\langle s,r\mid s^2=r^3=1\rangle$, let
$L=B_4/\langle z_4\rangle$, and let $P=\Gamma_{0,5}\triangleleft L$
be its pure subgroup, of index $24$.  Write $a=\sigma _1$,
$b=\sigma _2$, $c=\sigma _3$.  The geometric CM pullback, with its
ordinary homology basis fixed as in (FCE.4)--(FCE.6), identifies the geometric map
$\varphi$ of E.260 and the map $\rho_+$ of A.4.4 as isomorphisms
onto the same subgroup
\[
 G_3=q^{-1}(\langle ab\rangle)\subset L.                 \tag{MC.1}
\]
There are an actual element $g\in G_3$ and $\epsilon\in\{0,1\}$
such that
\[
 \varphi=\operatorname{Ad}(g)\,\rho_+\kappa^\epsilon,
 \qquad \kappa(s)=s,\quad \kappa(r)=r^{-1}.              \tag{MC.2}
\]
Here $g$ is retained as the inner marking supplied by the actual
automorphism $\rho_+^{-1}\varphi$; it is not set equal to $1$.

For every finite-index normal subgroup $N\triangleleft P$, put
$R=P/N$ and let $\mu$ be braid mirroring, $\mu(\sigma_i)=\sigma_i^{-1}$.
Choose one representative $t$ of each coset in $L/P$ and define
\[
 N^*=\bigcap_{t\in L/P}\bigl(tNt^{-1}\cap\mu(tNt^{-1})\bigr).
                                                               \tag{MC.3}
\]
Then $N^*\triangleleft L$, $\mu(N^*)=N^*$, and the actual diagonal
map
\[
 P/N^*\longrightarrow R^{48},\qquad
 pN^*\longmapsto
 \bigl(t^{-1}\mu^\eta(p)t\,N\bigr)_{t\in L/P,\ \eta=0,1}
                                                               \tag{MC.4}
\]
is injective.  In particular
\[
 |L/N^*|\leq 24|R|^{48}.                                \tag{MC.5}
\]
All inner markings and the mirror marking act on this finite
refinement by specified automorphisms.  Thus the finite equations
and their source frame labels for $\varphi$ are transported exactly
to those for $\rho_+$ on $L/N^*$; the projection to the $t=1,\eta=0$
factor in (MC.4) recovers the originally calculated quotient $R$.
This supplies the finite marking comparison for the successive-prime
factors as well as for the full cubic detector.  It supplies no
pentagon exclusion without the corresponding finite calculation.
\end{result}

\begin{proof}
First we prove the exact automorphism assertion used in (MC.2).
Every finite subgroup of $C_2*C_3$ fixes a vertex of its Bass--Serre
tree: the convex hull of one finite orbit is a finite tree, whose
centre is fixed, and the action has no edge inversions.  Its vertex
stabilizers are the conjugates of $C_2$ and $C_3$.  An automorphism
therefore sends $s$ to a conjugate of $s$ and $r$ to a conjugate of
$r$ or $r^{-1}$.  After one inner conjugation it fixes $s$ and sends
$r$ to $w r^\epsilon w^{-1}$, where here $\epsilon\in\{1,-1\}$.

Reduce $w$ in its double coset $\langle s\rangle w\langle r\rangle$.
Right multiplication by a power of $r$ does not change the image of
$r$.  Left multiplication by $s$ amounts to one further inner
conjugation, which still fixes $s$.  A nonempty reduced representative
$w$ consequently begins with an $r$-syllable and ends with an
$s$-syllable.  In an alternating word in $s$ and
$w r^{\pm1}w^{-1}$, every internal seam is then reduced: the seams
$w r^{\pm1}$ and $r^{\pm1}w^{-1}$ meet different factors, as do the
seams $w^{-1}s w$.  The same is true at either external occurrence
of $s$.  Every element of
$\langle s,w r w^{-1}\rangle$ involving the second factor has
free-product syllable length at least $2|w|+1>1$.
It cannot equal $r$.  Surjectivity of the automorphism forces $w$
to be empty.  We have proved
\[
 \operatorname{Aut}(C_2*C_3)
   =\operatorname{Inn}(C_2*C_3)\rtimes\langle\kappa\rangle,
 \qquad \operatorname{Out}(C_2*C_3)\cong C_2.             \tag{MC.6}
\]
The two outer classes are different because they act differently
on the quotient $C_3$.  Apply (MC.6) to the actual automorphism
$\rho_+^{-1}\varphi$ furnished by (MC.1).  Its inner element has
an actual image in $G_3$, giving (MC.2).

We next express the nontrivial outer class by ambient braid maps,
so that it is available on the finite target.  Braid mirroring is
an automorphism of $L$ because it sends $z_4$ to $z_4^{-1}$.
It preserves $P$, since an inverse standard braid has the same
permutation as that braid.  It commutes with folding, and
\[
 \mu(ab)=a^{-1}b^{-1}=a^{-1}(ab)^{-1}a.
\]
Consequently $\iota=\operatorname{Ad}(a)\mu$ preserves $G_3$,
satisfies $\iota^2=1$, and induces inversion on its cyclic quotient
$G_3/\ker q=C_6$.  The automorphism
$\rho_+^{-1}\iota\rho_+$ therefore has the nontrivial outer class
in (MC.6), because inner automorphisms act trivially on $B^{\rm ab}=C_6$.
There is an actual $k\in G_3$ with
\[
 \iota\rho_+=\operatorname{Ad}(k)\rho_+\kappa.           \tag{MC.7}
\]
Combining (MC.2) and (MC.7) absorbs $k$ into a named inner marking:
\[
 \varphi=\operatorname{Ad}(g_0)\iota^\epsilon\rho_+
 \quad\text{for an actual }g_0\in G_3.                  \tag{MC.8}
\]
This is an equality of homomorphisms, including their based
peripheral elements, rather than an identification of unmarked
groups.

Because $N$ is normal in $P$, conjugation of $N$ by an element of
$L$ depends only on its coset in $L/P$.  The intersection
$N_0=\bigcap_t tNt^{-1}$ is therefore normal in $L$ and contained
in $P$.  Formula (MC.3) is exactly $N_0\cap\mu(N_0)$, so it is
normal in $L$ and invariant under $\mu$.  Each coordinate of (MC.4)
is a homomorphism, since $P$ is normal in $L$ and invariant under
$\mu$.  The intersection of their kernels is precisely (MC.3).
This proves injectivity and the bound (MC.5).  These arguments
also specify the quotient maps: no isomorphism between an unmarked
factor and its conjugate is assumed.

Let $\pi_*:L\to L/N^*$ be the quotient map.  Normality and mirror
invariance make both $\operatorname{Ad}(\pi_*(g_0))$ and the induced
$\overline\iota$ actual automorphisms.  Equation (MC.8) becomes
\[
 \pi_*\varphi=
 \operatorname{Ad}(\pi_*(g_0))\,
 \overline\iota^{\epsilon}\,\pi_*\rho_+.               \tag{MC.9}
\]
For a finite tuple, apply the same displayed automorphism to every
entry, including the prescribed torsion elements, peripheral words,
and conjugators.  It carries a word equation to its image word
equation, and its inverse gives the reverse transport.  It therefore
induces a bijection between the two exactly marked finite solution
sets.  Projection by (MC.4) gives each of the original conjugate
and mirror factor equations, with their literal transported words.

The parity label is transported with the frames.  Conjugation
preserves each scalar central sign in the CM matrices.  The outer
change $\kappa$ is inner-conjugate to the normalized complex
conjugation action.  On gentle parameters its effect is conjugation
by that actual element; the already established homomorphism
$\Delta$ consequently satisfies
$\Delta(c h c^{-1})=\Delta(c)+\Delta(h)-\Delta(c)=\Delta(h)$.
Thus the bijection just described preserves the parity-one source
condition.  In particular an empty parity-one solution set in the
fully specified factor equations has empty preimage under the
refinement and remains empty after (MC.9).  This is the exact
transfer used for a finite exclusion; it does not assert that the
factor equations computed so far have an empty solution set.
\end{proof}
